\chapter{The Selberg--Delange Method}\label{ch:sd}
\begin{center}\itshape
A square root of the zeta function sounds like nonsense.\\
Landau took one anyway, and the sums of two squares\\
came out thinned by exactly $\sqrt{\log x}$.
\end{center}

In 1908 Edmund Landau counted the integers that are sums of two squares. There are infinitely many of them, of course, and they have density zero; the question was how fast they thin out. Landau's answer was that their number up to $x$ is
\[
B(x)\sim b\,\frac x{\sqrt{\log x}},\qquad b=\frac1{\sqrt2}\prod_{p\equiv3\ (4)}\Bigl(1-\frac1{p^2}\Bigr)^{-1/2}=0.76422\ldots,
\]
and his proof contained a startling idea. The generating Dirichlet series of these integers behaves near $s=1$ like $\sqrt{\zeta(s)}$, which has no pole at $s=1$ but a \emph{branch point}, and Landau integrated around the branch cut.

Five years later, on 16 January 1913, a clerk at the Madras Port Trust sent a letter to G.~H.~Hardy in Cambridge. Among the dozens of formulas in it was the claim that $B(x)=b\int_2^x\frac{\dd t}{\sqrt{\log t}}+\theta(x)$ with $\theta(x)$ ``very small''. When Hardy asked what ``very small'' meant, Srinivasa Ramanujan wrote back: of order $\frac{\sqrt x}{\log x}$. That was false, as Daniel Shanks showed in 1964. But Hardy, relying on a computation by his student Geraldine Stanley that later turned out to be wrong, had written in his lectures on Ramanujan that Landau's approximation was the better one; Shanks showed that Ramanujan's is. Ramanujan was wrong in the right direction. Our own computer agrees: up to $10^7$ there are $1\,985\,459$ sums of two squares, Landau's formula gives $1\,903\,547$ and Ramanujan's gives $1\,969\,213$.

Forty years after Landau, in 1953 and 1954, the Indian mathematician L.~G.~Sathe published four long papers, more than a hundred and thirty pages in the \emph{Journal of the Indian Mathematical Society}, on a problem Hardy had raised: how many $n\le x$ have exactly $k$ prime factors, uniformly in $k$? Atle Selberg read them. In a five-page note printed in the same volume as Sathe's last two parts, he wrote that the results were very beautiful but the method complicated, and gave a new proof: study $\sum_{n\le x}z^{\omega(n)}$ for complex $z$, whose Dirichlet series is $\zeta(s)^z$ times a tame factor, and integrate around the branch cut of $\zeta(s)^z$ at $s=1$, exactly as Landau had done for $z=\frac12$. Hubert Delange turned the idea into a general theorem in 1959 and 1971, and the result is now called the \emph{Selberg--Delange method}.

Tome~10 (Chapter~11, equation~(11.1)) stated its conclusion for $\omega(n)$ and promised the proof here. The answer has a probabilistic face: the mean of $z^{\omega(n)}$ is the residue of a Poisson variable, corrected by an arithmetic factor and by $\frac1{\Gamma(z)}$, and that Gamma factor comes from the sharp cut-off $n\le x$ and from nothing else. Hankel's contour integral for $\frac1{\Gamma(z)}$ is the bridge.

\begin{reminder}[from Tomes~1, 4, 5 and~13]
\emph{Gamma} (Tome~4, Theorem~9.36). $\Gamma$ is meromorphic on $\C$ with poles at $0,-1,-2,\dots$, $\Gamma(z)\Gamma(1-z)=\frac\pi{\sin\pi z}$ for $z\notin\Z$, and $\frac1\Gamma$ is entire.
\emph{Zeros of $\zeta$} (Tome~5, Proposition~6.11 and Theorem~6.12). For $-1\le\sigma\le2$ and $|t|\ge2$, $\frac{\zeta'}\zeta(s)=\sum_{|\gamma-t|\le1}\frac1{s-\rho}+O(\log|t|)$, the sum having $O(\log|t|)$ terms; and $\zeta(\sigma+it)\ne0$ for $\sigma\ge1-\frac c{\log(|t|+2)}$, with an absolute $c>0$. Also $-\frac{\zeta'}\zeta(\sigma)\le\frac1{\sigma-1}+A$ for $1<\sigma\le2$.
\emph{Zeros of $L(s,\chi)$} (Tome~13, Lemma~2.2 and Theorem~2.3). For $\chi$ non-principal modulo $q$, $\mathcal L=\log(q(|t|+2))$ and $\frac78\le\sigma\le\frac{11}8$, $\frac{L'}L(s,\chi)=\sum_{|\rho-s_0|\le1/2}\frac1{s-\rho}+O(\mathcal L)$ with $s_0=\frac98+it$, the sum having $O(\mathcal L)$ terms; and $L(s,\chi)\ne0$ for $\sigma\ge1-\frac{c_1}{\mathcal L}$, except possibly for one real zero when $\chi$ is real.
\emph{Perron} (Tome~5, Lemma~7.1). For $y,c,T>0$ let $I(y,T)=\frac1{2\pi i}\int_{c-iT}^{c+iT}\frac{y^s}s\dd s$. Then $|I(y,T)-\mathbf 1_{y>1}|\le y^c\min(1,\frac1{\pi T|\log y|})$ for $y\ne1$, and $|I(1,T)-\frac12|\le\frac c{\pi T}$.
\emph{Two squares} (Tome~1, Theorem~7.3). $n$ is a sum of two squares if and only if every prime $q\equiv3\pmod4$ divides $n$ to an even power.
\end{reminder}

\section{Hankel's loop}

The whole method rests on one formula from the theory of the Gamma function. For $W>0$, let $\mathcal H_W$ be the path that comes from $-W$ along the lower edge of the negative real axis, goes once counterclockwise around the circle $|w|=1$, and returns to $-W$ along the upper edge; $\mathcal H=\mathcal H_\infty$ is \emph{Hankel's loop}. On it, $w^{-z}=e^{-z\log w}$ with the principal logarithm, $\arg w\in[-\pi,\pi]$, where $\arg w=-\pi$ on the lower edge and $\pi$ on the upper one.

\begin{lemma}[Hankel, 1864]\label{lem:hankel}
For every $z\in\C$,
\[
\frac1{2\pi i}\int_{\mathcal H}w^{-z}e^w\dd w=\frac1{\Gamma(z)}.
\]
Moreover, for $|z|\le R$ and $W\ge1$,
\[
\Bigl|\frac1{2\pi i}\int_{\mathcal H\setminus\mathcal H_W}w^{-z}e^w\dd w\Bigr|\le\frac{e^{\pi R}}\pi\int_W^\infty u^Re^{-u}\dd u.
\]
\end{lemma}
\begin{proof}
On the edges, $w=ue^{\pm i\pi}$ with $u\ge1$, and $|w^{-z}|=u^{-\operatorname{Re}z}e^{\pm\pi\operatorname{Im}z}\le e^{\pi R}u^R$ when $|z|\le R$, while $|e^w|=e^{-u}$. This gives the tail bound, and shows that the integral converges absolutely and uniformly for $|z|\le R$. The integrand is entire in $z$, so the integral is an entire function of $z$ (by Fubini and Morera's theorem, Tome~4, Chapter~9).

The circle $|w|=1$ can be replaced by $|w|=r$ for any $0<r\le1$: the difference is the integral over the boundary of a region where $w^{-z}e^w$ is holomorphic (the annulus slit along the negative axis), which vanishes by Cauchy's theorem. Let $\operatorname{Re}z<1$. On $|w|=r$ the integrand is at most $r^{-\operatorname{Re}z}e^{\pi|\operatorname{Im}z|}e^r$, and the circle has length $2\pi r$, so its contribution tends to $0$ as $r\to0$. On the lower edge, travelled from $-\infty$ to $0$, $w=ue^{-i\pi}$ and $w^{-z}=u^{-z}e^{i\pi z}$; on the upper edge, travelled from $0$ to $-\infty$, $w^{-z}=u^{-z}e^{-i\pi z}$. With $\dd w=-\dd u$,
\[
\int_{\mathcal H}w^{-z}e^w\dd w=e^{i\pi z}\int_0^\infty u^{-z}e^{-u}\dd u-e^{-i\pi z}\int_0^\infty u^{-z}e^{-u}\dd u=2i\sin(\pi z)\,\Gamma(1-z).
\]
So for $\operatorname{Re}z<1$ and $z\notin\Z$ the left side, divided by $2\pi i$, is $\frac{\sin\pi z}\pi\Gamma(1-z)=\frac1{\Gamma(z)}$ by the reflection formula. Both sides are entire, and they agree on an open set, so they agree everywhere.
\end{proof}

For $z=k$ a positive integer, $w^{-k}e^w$ is single-valued, the two edges cancel, and the lemma says that the residue of $e^ww^{-k}$ at $0$ is $\frac1{(k-1)!}$; for $z=0,-1,-2,\dots$ the integrand is entire and the integral is $0=\frac1{\Gamma(z)}$. Hankel's formula interpolates between these.

\begin{think}
Why should a sharp cut-off $n\le x$ produce $\frac1{\Gamma(z)}$? Compute $\frac1{2\pi i}\int_{c-i\infty}^{c+i\infty}(s-1)^{-z}\frac{x^s}s\dd s$ heuristically by the substitution $s=1+\frac w{\log x}$. What would the weight $e^{-n/x}$ instead of $\mathbf 1_{n\le x}$ give?
\end{think}

\section{Powers of zeta}

Let $c$ and $c_1$ be the absolute constants of Tome~5, Theorem~6.12, and Tome~13, Theorem~2.3. Put $c_2=\min(c,\frac{c_1}3)$ and $c_*=\min(\frac{c_2}2,\frac18)$, write $\tau=|t|+3$, and define the closed region and the open region
\[
\mathcal D=\Bigl\{s:\sigma\ge1-\frac{c_*}{\log\tau}\Bigr\}\subset U=\Bigl\{s:\sigma>1-\frac{2c_*}{\log\tau}\Bigr\}.
\]
Both lie in the half-plane $\sigma>\frac34$, because $\frac{2c_*}{\log3}\le\frac14\cdot\frac1{\log 3}<\frac14$. The region $U$ lies to the right of a curve, so it is simply connected.

Let $\chi_4$ be the non-principal character modulo $4$. We need three facts about $U$.
\begin{itemize}
\item $\zeta$ has no zero in $U$: there $\sigma>1-\frac{2c_*}{\log(|t|+3)}\ge1-\frac c{\log(|t|+2)}$.
\item $L(s,\chi_4)$ has no zero in $U$. Indeed $\log(4(|t|+2))\le3\log\tau$, so $U$ lies in the region of Tome~13, Theorem~2.3, with $q=4$; and the possible exceptional zero would be real, while $L(\sigma,\chi_4)=1-3^{-\sigma}+5^{-\sigma}-\cdots>0$ for real $\sigma>0$ (an alternating series with decreasing terms).
\item The function $Z(s)=(s-1)\zeta(s)$ is entire and has no zero in $U$, since $Z(1)=1$.
\end{itemize}
On the simply connected $U$, the non-vanishing holomorphic functions $Z$ and $L(\cdot,\chi_4)$ have holomorphic logarithms (Tome~4, Chapter~9), which we normalise to be real on $U\cap\R$. On the slit region $U^-=U\setminus\left]-\infty,1\right]$ put
\[
\log\zeta(s)=\log Z(s)-\log(s-1),\qquad \zeta(s)^z=e^{z\log\zeta(s)}=Z(s)^z(s-1)^{-z},
\]
with the principal $\log(s-1)$. This is a logarithm of $\zeta$, real on $\left]1,\infty\right[$, so for $\sigma>1$ it is the logarithm $\sum_n\frac{\Lambda(n)}{n^s\log n}=-\sum_p\log(1-p^{-s})$ of Tome~5, since two logarithms of the same function differ by a constant multiple of $2\pi i$ and these agree on the real axis. In particular, for $\sigma>1$, $\zeta(s)^z=\prod_p(1-p^{-s})^{-z}$ with principal powers, and $|\log\zeta(s)|\le\log\zeta(\sigma)$. Likewise $\log L(s,\chi_4)=-\sum_p\log(1-\chi_4(p)p^{-s})$ for $\sigma>1$, of modulus at most $\log\zeta(\sigma)$.

The growth of these logarithms is what the method needs. Write $f$ for either $\zeta$ or $L(\cdot,\chi_4)$.

\begin{lemma}\label{lem:logderivK}
There is an absolute constant $A_1$ such that $|\frac{f'}f(s)|\le A_1\log\tau$ for $s\in\mathcal D$ with $|t|\ge12$.
\end{lemma}
\begin{proof}
If $\sigma\ge2$, $|\frac{f'}f(s)|\le\sum_n\Lambda(n)n^{-2}$. Let $\sigma_1=1+\frac1{\log\tau}$ and $s_1=\sigma_1+it$; since $\log\tau\ge\log15>\frac83$, $\sigma_1<\frac{11}8$. For $\sigma\ge\sigma_1$, $|\frac{f'}f(s)|\le\sum_n\Lambda(n)n^{-\sigma}=-\frac{\zeta'}\zeta(\sigma)\le-\frac{\zeta'}\zeta(\sigma_1)\le\log\tau+A$.

Let $1-\frac{c_*}{\log\tau}\le\sigma\le\sigma_1$; then $\sigma\ge\frac78$. By the reminders, at $s$ and at $s_1$,
\[
\frac{f'}f(s)=\sum_{\rho\in\mathcal Z}\frac1{s-\rho}+O(\log\tau),\qquad\frac{f'}f(s_1)=\sum_{\rho\in\mathcal Z}\frac1{s_1-\rho}+O(\log\tau),
\]
where $\mathcal Z$ is a set of $O(\log\tau)$ zeros $\rho=\beta+i\gamma$ of $f$ with $|\gamma-t|\le1$ (the same set at $s$ and $s_1$, because it only depends on $t$). For such $\rho$, $\log(|\gamma|+2)\le\log\tau$ and, for $f=L(\cdot,\chi_4)$, $\log(4(|\gamma|+2))\le3\log\tau$; so the zero-free regions give $\beta\le1-\frac{c_2}{\log\tau}$.

\emph{Zeros seen from $s_1$.} Taking real parts at $s_1$, and using $|\frac{f'}f(s_1)|\le\log\tau+A$,
\[
\sum_{\rho\in\mathcal Z}\frac{\sigma_1-\beta}{|s_1-\rho|^2}=\operatorname{Re}\sum_{\rho\in\mathcal Z}\frac1{s_1-\rho}=O(\log\tau).
\]
Each $\sigma_1-\beta\ge\sigma_1-1=\frac1{\log\tau}$, so $\sum_{\rho\in\mathcal Z}|s_1-\rho|^{-2}=O(\log^2\tau)$.

\emph{From $s_1$ to $s$.} Since $\sigma\ge1-\frac{c_*}{\log\tau}\ge1-\frac{c_2}{2\log\tau}$, we have $\sigma-\beta\ge\frac{c_2}{2\log\tau}$, while $\sigma_1-\sigma\le\frac{1+c_*}{\log\tau}\le\frac2{\log\tau}$. So $\sigma_1-\beta\le(1+\frac4{c_2})(\sigma-\beta)$, and since $s$ and $s_1$ have the same imaginary part, $|s_1-\rho|\le(1+\frac4{c_2})|s-\rho|$. Subtracting the two expansions,
\[
\begin{aligned}\Bigl|\frac{f'}f(s)-\frac{f'}f(s_1)\Bigr|&\le\sum_{\rho\in\mathcal Z}\frac{|s_1-s|}{|s-\rho||s_1-\rho|}+O(\log\tau)\\&\le\frac2{\log\tau}\Bigl(1+\frac4{c_2}\Bigr)\sum_{\rho\in\mathcal Z}\frac1{|s_1-\rho|^2}+O(\log\tau)=O(\log\tau).\end{aligned}\qedhere
\]
\end{proof}

\begin{lemma}\label{lem:logzeta}
There is an absolute constant $A_2$ such that for $s\in\mathcal D$:
\begin{enumerate}[label=(\alph*)]
\item if $|t|\ge12$, then $|\log\zeta(s)|\le\log\log\tau+A_2$ and $|\log L(s,\chi_4)|\le\log\log\tau+A_2$;
\item if $|t|\le12$, then $|\log L(s,\chi_4)|\le A_2$, and $|\log Z(s)|\le A_2$ when moreover $\sigma\le2$.
\end{enumerate}
Consequently, for $|z|\le R$ and $|t|\ge12$, $|\zeta(s)^z|\le e^{RA_2}(\log\tau)^R$.
\end{lemma}
\begin{proof}
(a) For $\sigma\ge\sigma_1=1+\frac1{\log\tau}$, $|\log f(s)|\le\log\zeta(\sigma_1)\le\log\frac{\sigma_1}{\sigma_1-1}\le\log(2\log\tau)$, using $\zeta(\sigma)\le1+\int_1^\infty u^{-\sigma}\dd u=\frac\sigma{\sigma-1}$. For $1-\frac{c_*}{\log\tau}\le\sigma<\sigma_1$, $\log f(s)=\log f(s_1)-\int_\sigma^{\sigma_1}\frac{f'}f(u+it)\dd u$, and the integral is at most $\frac2{\log\tau}A_1\log\tau=2A_1$ by Lemma~\ref{lem:logderivK}. Take $A_2\ge\log2+2A_1$.

(b) The set $\{s\in\mathcal D:|t|\le12,\ \sigma\le2\}$ is compact and contained in $U$, where $\log Z$ and $\log L(\cdot,\chi_4)$ are continuous, hence bounded. For $\sigma\ge2$, $|\log L(s,\chi_4)|\le\log\zeta(2)$.

Finally $|\zeta(s)^z|=e^{\operatorname{Re}(z\log\zeta(s))}\le e^{R|\log\zeta(s)|}$.
\end{proof}

\section{The theorem}

For an integer $K\ge1$, let $\tau_K(n)$ be the number of ordered factorisations $n=d_1\cdots d_K$, so that $\sum_n\tau_K(n)n^{-s}=\zeta(s)^K$ for $\sigma>1$; it is multiplicative, and $\tau_K(p^a)$ is the number of $\mathbf u\in\N^K$ with $u_1+\dots+u_K=a$.

\begin{lemma}\label{lem:taupower}
For $K\ge1$, $n\ge1$ and $y\ge1$: $\tau_K(n)^2\le\tau_{K^2}(n)$, and $\sum_{n\le y}\tau_K(n)\le y(1+\log y)^K$.
\end{lemma}
\begin{proof}
Both sides of the first inequality are multiplicative, so let $n=p^a$. To a pair $(\mathbf u,\mathbf v)$ with $\sum u_i=\sum v_j=a$ we attach a $K\times K$ matrix of non-negative integers with row sums $\mathbf u$ and column sums $\mathbf v$; it exists by induction on $a$ (if $a\ge1$, pick $u_i\ge1$ and $v_j\ge1$, put a $1$ in position $(i,j)$, and apply the induction to $(\mathbf u-e_i,\mathbf v-e_j)$). Its row and column sums recover $(\mathbf u,\mathbf v)$, so the map is injective into the matrices with total $a$, which number $\tau_{K^2}(p^a)$.

For the second, $\sum_{n\le y}\tau_K(n)\le y\sum_{n\le y}\frac{\tau_K(n)}n\le y\bigl(\sum_{d\le y}\frac1d\bigr)^K$, since every factorisation $n=d_1\cdots d_K$ with $n\le y$ has all $d_i\le y$; and $\sum_{d\le y}\frac1d\le1+\log y$.
\end{proof}

\begin{theorem}[Selberg--Delange]\label{thm:sd}
Let $R\ge1$, let $K$ be the least integer $\ge R$, and let $z\in\C$ with $|z|\le R$. Let $(a_n)$ be complex numbers with $|a_n|\le\tau_K(n)$, and $F(s)=\sum_na_nn^{-s}$ for $\sigma>1$. Suppose that $F(s)=\zeta(s)^zG(s)$ for $\sigma>1$, where $G$ is holomorphic on $U$ and $|G(s)|\le M(\log\tau)^R$ on $\mathcal D$. Then for $x\ge3$,
\[
\sum_{n\le x}a_n=x(\log x)^{z-1}\Bigl(\frac{G(1)}{\Gamma(z)}+O_R\Bigl(\frac{M+1}{\log x}\Bigr)\Bigr),
\]
where the implied constant depends only on $R$.
\end{theorem}

Before the proof, a word on its shape. The sum is a Perron integral of $F(s)\frac{x^s}s$ along a vertical line. Pushing the line to the left of $1$ is impossible because of the branch cut of $\zeta(s)^z$, so the line wraps around the cut (Figure~\ref{fig:hankel}). Everything away from $s=1$ is small because $x^s$ is small there; the loop around the cut, rescaled by $\log x$, becomes Hankel's loop.

\begin{figure}[ht]
\centering\fig{ch02-hankel}
\caption{The contour of the proof. The Perron segment on $\sigma=c$ is replaced by the three sides of a rectangle, with a loop around the cut $[\sigma_T,1]$ of $\zeta(s)^z$. After the change of variables $s=1+\frac w{\log x}$ the loop becomes the truncated Hankel loop $\mathcal H_W$.}
\label{fig:hankel}
\end{figure}

\begin{proof}
Write $L=\log x$. All constants $C_1,C_2,\dots$ and all $O$'s depend only on $R$. Put $B=4(K^2+R+3)$ and $T=L^B$. We first assume $x\ge x_0(R)$, where $x_0$ is so large that $T\le x$, $T\ge12$, $\log(T+3)\ge4$ and $\frac1L<\frac{c_*}{\log(T+3)}$; the small values of $x$ are treated at the end.

\emph{Step 1: Perron's formula.} Let $c=1+\frac1L$. Since $\sum|a_n|n^{-c}\le\zeta(c)^K<\infty$ we may integrate term by term, and $\frac1{2\pi i}\int_{c-iT}^{c+iT}F(s)\frac{x^s}s\dd s=\sum_na_nI(\frac xn,T)$. The difference $E_1$ between this and $\sum_{n\le x}a_n$ is at most $\sum_n|a_n||\mathbf 1_{n\le x}-I(\frac xn,T)|$. Let $u=T^{-1/2}$. For $|\log\frac xn|\ge u$, Perron's lemma bounds the term by $|a_n|(\frac xn)^c\frac1{\pi Tu}$, and these terms add up to at most
\[
\frac{x^c}{\sqrt T}\sum_n\frac{|a_n|}{n^c}\le\frac{e\,x}{\sqrt T}\zeta(c)^K\le\frac{e\,x(2L)^K}{\sqrt T},
\]
since $\zeta(c)\le\frac c{c-1}\le2L$. For $|\log\frac xn|<u$, the term is at most $2|a_n|$ (for $n\ne x$, $(\frac xn)^c\le e^{2u}<2$; for $n=x$, $|1-I(1,T)|\le\frac12+\frac c{\pi T}<1$). These $n$ lie in an interval of length $x(e^u-e^{-u})\le3xu$, so there are at most $3xu+1\le4xu$ of them (as $xu\ge\sqrt x\ge1$). By Cauchy--Schwarz and Lemma~\ref{lem:taupower},
\[
\sum_{|\log\frac xn|<u}|a_n|\le(4xu)^{1/2}\Bigl(\sum_{n\le2x}\tau_K(n)^2\Bigr)^{1/2}\le(4xu)^{1/2}\bigl(2x(2L)^{K^2}\bigr)^{1/2}\le3xT^{-1/4}(2L)^{K^2/2}.
\]
Altogether $|E_1|\le9\cdot2^{K^2}xT^{-1/4}L^{K^2}=9\cdot2^{K^2}xL^{-R-3}$.

\emph{Step 2: the contour.} Let $\delta=\frac{c_*}{\log(T+3)}$ and $\sigma_T=1-\delta$. For $|t|\le T$, $1-\frac{c_*}{\log\tau}\le\sigma_T$, so the rectangle $[\sigma_T,c]\times[-T,T]$ lies in $\mathcal D$, where $\zeta(s)^zG(s)\frac{x^s}s$ is holomorphic except on the cut $[\sigma_T,1]$; it coincides with $F(s)\frac{x^s}s$ on $\sigma=c$. By Cauchy's theorem, the integral over $[c-iT,c+iT]$ equals the sum of the integrals over the two horizontal sides $J_\pm$, the two vertical half-sides $V_\pm$ on $\sigma=\sigma_T$, and the loop $\Lambda$ around the cut made of its two edges and the circle $|s-1|=\frac1L$, all as in Figure~\ref{fig:hankel}. (Apply Cauchy's theorem to the rectangle minus a thin neighbourhood of the cut and let it shrink; the integrand is continuous up to the cut from each side.)

\emph{Step 3: the far part.} On $J_\pm$, $|t|=T\ge12$, so by Lemma~\ref{lem:logzeta} $|\zeta(s)^zG(s)|\le e^{RA_2}M(\log(T+3))^{2R}$, while $|\frac{x^s}s|\le\frac{ex}T$ and the length is at most $1$. So $|J_\pm|\le C_1Mx(\log(T+3))^{2R}T^{-1}$.

On $V_\pm$, $|x^s|=x^{\sigma_T}=xe^{-c_*L/\log(T+3)}$. Where $|t|\ge12$, the bound of $\zeta^zG$ is the same as on $J_\pm$. Where $|t|\le12$, $|\zeta(s)^z|=|Z(s)^z||(s-1)^{-z}|\le e^{RA_2}e^{\pi R}\max(|s-1|,|s-1|^{-1})^R\le C_2(\log(T+3))^R$, because $\delta\le|s-1|\le13$, and $|G(s)|\le M(\log 15)^R$. Since $|s|\ge\max(\frac12,|t|)$, $\int_{-T}^T\frac{\dd t}{|s|}\le4+2\log T$. So
\[
|V_+|+|V_-|\le C_3Mx(\log(T+3))^{2R+1}e^{-c_*L/\log(T+3)}.
\]
Now $\log(T+3)\le2B\log L$, so $e^{-c_*L/\log(T+3)}\le e^{-c_*L/(2B\log L)}$, which is $O(L^{-A})$ for every fixed $A$. Hence $|J_\pm|+|V_\pm|\le C_4MxL^{-R-3}$.

\emph{Step 4: the loop.} Write $\zeta(s)^zG(s)\frac1s=(s-1)^{-z}H(s)$ with $H(s)=Z(s)^z\frac{G(s)}s$. Let $r_0=\frac{c_*}2$. The disc $|s-1|\le r_0$ lies in $\mathcal D$ and in the compact set of Lemma~\ref{lem:logzeta}(b), so there $|H(s)|\le2e^{RA_2}(\log 4)^RM\le C_5M$; by Cauchy's estimate $|H'(s)|\le\frac{2C_5M}{r_0}$ on $|s-1|\le\frac{r_0}2$, so
\[
|H(s)-H(1)|\le\frac{2C_5M}{r_0}|s-1|\qquad(|s-1|\le\tfrac{r_0}2),\qquad H(1)=G(1),
\]
since $Z(1)=1$ and $\log Z(1)=0$. The loop $\Lambda$ lies in $|s-1|\le\delta\le\frac{r_0}2$, because $\log(T+3)\ge4$. Substitute $s=1+\frac wL$. The loop becomes $\mathcal H_W$ with $W=\delta L$, $x^s=xe^w$, $\dd s=\frac{\dd w}L$, and $(s-1)^{-z}=L^zw^{-z}$ (the principal arguments of $\frac wL$ and $w$ coincide). So
\[
\frac1{2\pi i}\int_\Lambda=xL^{z-1}\,\frac1{2\pi i}\int_{\mathcal H_W}w^{-z}e^wH\Bigl(1+\frac wL\Bigr)\dd w=xL^{z-1}\bigl(G(1)\,\mathcal I_1+\mathcal I_2\bigr),
\]
with
\[
\mathcal I_1=\frac1{2\pi i}\int_{\mathcal H_W}w^{-z}e^w\dd w,\qquad|\mathcal I_2|\le\frac{C_5M}{\pi r_0L}\int_{\mathcal H_W}|w||w^{-z}||e^w||\dd w|.
\] On the circle $|w|=1$ the integrand of $\mathcal I_2$ is at most $e^{\pi R+1}$; on the edges it is at most $e^{\pi R}u^{R+1}e^{-u}$. So $|\mathcal I_2|\le C_6\frac ML$. By Lemma~\ref{lem:hankel}, $\mathcal I_1=\frac1{\Gamma(z)}+O(W^Re^{-W})$, and $W=\frac{c_*L}{\log(T+3)}\ge\frac{c_*L}{2B\log L}$, so $W^Re^{-W}=O(L^{-1})$. Also $|G(1)|\le M(\log3)^R\le2^RM$.

\emph{Step 5: conclusion for large $x$.} Collecting,
\[
\sum_{n\le x}a_n=xL^{z-1}\Bigl(\frac{G(1)}{\Gamma(z)}+O\Bigl(\frac ML\Bigr)\Bigr)+O\bigl((M+1)xL^{-R-3}\bigr),
\]
and $|xL^{z-1}|=xL^{\operatorname{Re}z-1}\ge xL^{-R-1}$, so the last error is $O((M+1)|xL^{z-1}|L^{-2})$.

\emph{Small $x$.} For $3\le x<x_0(R)$, $L$ lies in the fixed interval $[\log3,\log x_0]$. The left side is at most $\sum_{n\le x}\tau_K(n)\le x(1+L)^K$ by Lemma~\ref{lem:taupower}, and $|xL^{z-1}\frac{G(1)}{\Gamma(z)}|\le C_7Mx$ since $\frac1\Gamma$ is bounded on $|z|\le R$. So the difference between the two sides is $O((M+1)x)$, which is $O((M+1)|xL^{z-1}|L^{-1})$ because $|L^{z-1}|L^{-1}\ge L^{-R-2}$ is bounded below.
\end{proof}

\begin{remark}
Delange's theorem is sharper: under the same hypotheses there is an asymptotic expansion $\sum_{n\le x}a_n=xL^{z-1}\bigl(\sum_{j<J}\frac{\lambda_j(z)}{L^j}+O(L^{-J})\bigr)$, with $\lambda_j(z)=\frac1{\Gamma(z-j)}\frac1{j!}\frac{\dd^j}{\dd s^j}\bigl[(s-1)^z\zeta(s)^z\frac{G(s)}s\bigr]_{s=1}$, obtained by expanding $H$ to order $J$ in Step~4 (Tenenbaum, Chapter~II.5). The term $j=1$ explains the slow convergence in Figure~\ref{fig:sdratio}: for $z=3$ the ratio is still $1.45$ at $x=10^7$.
\end{remark}

\section{Prime factors: the promise of Tome~10}

\begin{lemma}\label{lem:eulerprod}
Let $R\ge1$ and $|z|\le R$. For $\sigma>1$, $\sum_nz^{\omega(n)}n^{-s}=\zeta(s)^zG_z(s)$, where
\[
G_z(s)=\prod_p\Bigl(1+\frac z{p^s-1}\Bigr)\Bigl(1-\frac1{p^s}\Bigr)^z
\]
converges absolutely and uniformly on $\sigma\ge\frac34$, is holomorphic on $\sigma>\frac34$, and satisfies $|G_z(s)|\le M(R)$ there. At $s=1$,
\[
G_z(1)=\prod_p\Bigl(1+\frac{z-1}p\Bigr)\Bigl(1-\frac1p\Bigr)^{z-1}.
\]
\end{lemma}
\begin{proof}
Since $|z^{\omega(n)}|\le K^{\omega(n)}\le\tau_K(n)$, the series converges absolutely for $\sigma>1$ and, $z^{\omega(n)}$ being multiplicative, it has the Euler product $\prod_p(1+z\sum_{k\ge1}p^{-ks})=\prod_p\phi(p^{-s})$, where $\phi(u)=1+\frac{zu}{1-u}=(1-u+zu)(1-u)^{-1}$. Multiplying by $\prod_p(1-p^{-s})^z=\zeta(s)^{-z}$ gives the identity, with factors $g_p(s)=\psi(p^{-s})$ where
\[
\psi(u)=(1+(z-1)u)(1-u)^{z-1},\qquad|u|<1.
\]
Now $\psi(0)=1$ and $\psi'(0)=(z-1)-(z-1)=0$, so $\frac{\psi(u)-1}{u^2}$ is holomorphic on $|u|<1$. On $|u|=0.8$, $|1-u|\ge0.2$ and $|\log(1-u)|\le\log5+\frac\pi2$, so $|\psi(u)-1|\le1+1.8(R+1)e^{(R+1)(\log5+\pi/2)}=:C_R$; by the maximum principle, $|\psi(u)-1|\le\frac{C_R}{0.64}|u|^2$ for $|u|\le0.8$. For $\sigma\ge\frac34$, $|p^{-s}|\le2^{-3/4}<0.6$, so $|g_p(s)-1|\le2C_Rp^{-3/2}$. Since $\sum_pp^{-3/2}<2$, the product converges absolutely and uniformly, its limit is holomorphic, and $|G_z(s)|\le\prod_p(1+2C_Rp^{-3/2})\le e^{4C_R}$. Finally $(1+\frac z{p-1})(1-\frac1p)^z=\frac{p-1+z}p(1-\frac1p)^{z-1}$.
\end{proof}

\begin{corollary}[Selberg, 1954]\label{cor:omegasd}
Uniformly for $|z|\le R$ and $x\ge3$,
\[
\frac1x\sum_{n\le x}z^{\omega(n)}=(\log x)^{z-1}\Bigl(\frac{F(z)}{\Gamma(z)}+O_R\Bigl(\frac1{\log x}\Bigr)\Bigr),\qquad F(z)=\prod_p\Bigl(1+\frac{z-1}p\Bigr)\Bigl(1-\frac1p\Bigr)^{z-1}.
\]
This is equation~(11.1) of Tome~10.
\end{corollary}
\begin{proof}
Apply Theorem~\ref{thm:sd} with $a_n=z^{\omega(n)}$ and $G=G_z$: since $U\subset\{\sigma>\frac34\}$, Lemma~\ref{lem:eulerprod} gives the hypotheses with $M=M(R)$.
\end{proof}

In the probabilistic language of Tome~10: for $n$ uniform in $[1,x]$, $\mathbb E_xz^{\omega(n)}=e^{(z-1)\log\log x}\bigl(\frac{F(z)}{\Gamma(z)}+o(1)\bigr)$, the generating function of a Poisson variable of parameter $\log\log x$, times a residual factor. Chapter~3 extracts from it the sharp form of the Erd\H{o}s--Kac theorem and the Sathe--Selberg formula for $\#\{n\le x:\omega(n)=k\}$.

\begin{figure}[ht]
\centering\fig{ch02-ratio}
\caption{The ratio of $\sum_{n\le x}z^{\omega(n)}$ to the main term $x(\log x)^{z-1}\frac{F(z)}{\Gamma(z)}$ of Corollary~\ref{cor:omegasd}, for $z=\frac12$, $2$ and $3$, with $x$ from $10^3$ to $10^7$. The ratios tend to $1$, as slowly as $1+\frac{\text{const}}{\log x}$ predicts: at $10^7$ they are $1.034$, $1.080$ and $1.454$. For $z=2$, $F(2)=\prod_p(1-\frac1{p^2})=\frac6{\pi^2}$.}
\label{fig:sdratio}
\end{figure}

\section{Landau's sums of two squares}

\begin{corollary}[Landau, 1908]\label{cor:landau}
The number $B(x)$ of integers $n\le x$ that are sums of two squares satisfies
\[
B(x)=\frac{b\,x}{\sqrt{\log x}}\Bigl(1+O\Bigl(\frac1{\log x}\Bigr)\Bigr),\qquad b=\frac1{\sqrt2}\prod_{p\equiv3\ (4)}\Bigl(1-\frac1{p^2}\Bigr)^{-1/2}=0.764223\ldots
\]
\end{corollary}
\begin{proof}
Let $a_n=1$ if $n$ is a sum of two squares and $0$ otherwise. By Tome~1, Theorem~7.3, $a_n$ is multiplicative, and for $\sigma>1$
\[
F(s)=\sum_na_nn^{-s}=(1-2^{-s})^{-1}\prod_{p\equiv1\ (4)}(1-p^{-s})^{-1}\prod_{p\equiv3\ (4)}(1-p^{-2s})^{-1}.
\]
Let $G(s)=L(s,\chi_4)^{1/2}P(s)$ with $L^{1/2}=e^{\frac12\log L}$ and
\[
P(s)=(1-2^{-s})^{-1/2}\prod_{p\equiv3\ (4)}(1-p^{-2s})^{-1/2},
\]
principal powers; $P$ is holomorphic and $|P|\le(1-2^{-3/4})^{-1/2}\zeta(\frac32)^{1/2}$ on $\sigma>\frac34$. For $\sigma>1$, compare logarithms prime by prime in $\log F$ and $\frac12\log\zeta+\log G$: at $p=2$ both are $-\log(1-2^{-s})$; at $p\equiv1\pmod4$ both are $-\log(1-p^{-s})$; at $p\equiv3\pmod4$ the right side is $-\frac12\log(1-p^{-s})-\frac12\log(1+p^{-s})-\frac12\log(1-p^{-2s})=-\log(1-p^{-2s})$. So $F=\zeta^{1/2}G$.

Take $z=\frac12$, $R=K=1$ (so $|a_n|\le1=\tau_1(n)$). By Lemma~\ref{lem:logzeta}, $|L(s,\chi_4)^{1/2}|\le e^{A_2/2}(\log\tau)^{1/2}$ on $\mathcal D$ (for $|t|\le12$ use part~(b)), so $|G|\le M\log\tau$ for an absolute $M$. Theorem~\ref{thm:sd} gives $B(x)=x(\log x)^{-1/2}(\frac{G(1)}{\Gamma(1/2)}+O(\frac1{\log x}))$. Finally $L(1,\chi_4)=\frac\pi4$ (Leibniz; Tome~5, Example~3.11) and $\Gamma(\frac12)=\sqrt\pi$, so
\[
\frac{G(1)}{\sqrt\pi}=\frac1{\sqrt\pi}\Bigl(\frac\pi4\Bigr)^{1/2}\sqrt2\prod_{p\equiv3\ (4)}\Bigl(1-\frac1{p^2}\Bigr)^{-1/2}=b.\qedhere
\]
\end{proof}

\begin{beauty}
The constant $b$ holds three pieces of the integers in one formula. The $\frac1{\sqrt\pi}$ from $\Gamma(\frac12)$ is the cut-off, the same as for the cycles of a random permutation. The $\sqrt{\pi/4}$ is Leibniz's series, the value at $1$ of the $L$-function that tells apart the primes $1$ and $3$ modulo~$4$. And the product over $p\equiv3\pmod4$ counts the price of those primes having to come in pairs.
\end{beauty}

Ramanujan's version replaces $\frac x{\sqrt{\log x}}$ by $\int_2^x\frac{\dd t}{\sqrt{\log t}}=\frac x{\sqrt{\log x}}(1+\frac1{2\log x}+\cdots)$, which builds in a second-order term of Delange's expansion. Shanks computed the true second-order term. It is not the one the integral produces, so Ramanujan's claim about $\theta(x)$ is false; but it is closer to the integral's than to Landau's, which has none, so the integral is the better approximation for all large $x$. Exercise~\ref{ex:onlythree} counts the opposite set: integers all of whose prime factors are $3$ modulo~$4$.

\section{Exercises}

\begin{exercise}[$\star$]\label{ex:hankelint}
Check Lemma~\ref{lem:hankel} directly for $z=1,2,3$ by the residue theorem, and for $z=0,-1$ by Cauchy's theorem. Use it to compute $\frac1{2\pi i}\int_{\mathcal H}w^{-1/2}e^w\dd w$.
\end{exercise}

\begin{exercise}[$\star$]\label{ex:twoomega}
Deduce from Corollary~\ref{cor:omegasd} that $\sum_{n\le x}2^{\omega(n)}=\frac6{\pi^2}x\log x+O(x)$. Then prove it without complex analysis, from $2^{\omega(n)}=\sum_{d\mid n}\mu(d)^2$.
\end{exercise}

\begin{exercise}[$\star\star$]\label{ex:minusone}
Show that $\sum_{n\le x}(-1)^{\omega(n)}=O(\frac x{(\log x)^3})$. Which value of the Gamma function is responsible? The notebook finds $\frac1x\sum_{n\le x}(-1)^{\omega(n)}=-0.000165$ at $x=10^7$.
\end{exercise}

\begin{exercise}[$\star\star$]\label{ex:bigomegasd}
Let $|z|\le\frac32$. Show that $\sum_{n\le x}z^{\Omega(n)}=x(\log x)^{z-1}\bigl(\frac{F_\Omega(z)}{\Gamma(z)}+O(\frac1{\log x})\bigr)$ with $F_\Omega(z)=\prod_p(1-\frac1p)^z(1-\frac zp)^{-1}$. What goes wrong at $z=2$?
\end{exercise}

\begin{exercise}[$\star\star$]\label{ex:onlythree}
Show that the number of $n\le x$ all of whose prime factors are $\equiv3\pmod4$ is $\frac{c_3x}{\sqrt{\log x}}(1+O(\frac1{\log x}))$, with $c_3=\frac{\sqrt2}\pi\prod_{p\equiv3\ (4)}(1-\frac1{p^2})^{-1/2}=0.48652\ldots$
\end{exercise}

\begin{exercise}[$\star\star\star$]\label{ex:landaupik}
(Landau, 1900.) Let $\pi_k(x)=\#\{n\le x:\omega(n)=k\}$. Show, for each fixed $k\ge1$, that
\[
\pi_k(x)=\frac x{\log x}\,\frac{(\log\log x)^{k-1}}{(k-1)!}\Bigl(1+O_k\Bigl(\frac1{\log\log x}\Bigr)\Bigr).
\]
\end{exercise}

\section{Hints}

\begin{hint}{ex:hankelint}
For $z=k\ge1$ the two edges cancel, since $w^{-k}$ is single-valued. For $z=\frac12$, $\Gamma(\frac12)=\sqrt\pi$.
\end{hint}
\begin{hint}{ex:twoomega}
$F(2)=\frac6{\pi^2}$ and $\Gamma(2)=1$. For the elementary proof, swap sums and use $\sum_{d\le y}\mu(d)^2=\frac6{\pi^2}y+O(\sqrt y)$.
\end{hint}
\begin{hint}{ex:minusone}
$\frac1{\Gamma(-1)}=0$.
\end{hint}
\begin{hint}{ex:bigomegasd}
The Euler factor is $(1-zp^{-s})^{-1}$; follow Lemma~\ref{lem:eulerprod} with $\psi(u)=(1-u)^z(1-zu)^{-1}$ on $|u|\le0.6$, and check that $|zu|\le0.9$.
\end{hint}
\begin{hint}{ex:onlythree}
Find $G$ with $F=\zeta^{1/2}G$; this time $L(s,\chi_4)$ appears with the exponent $-\frac12$.
\end{hint}
\begin{hint}{ex:landaupik}
Write
\[
\pi_k(x)=\frac1{2\pi i}\oint_{|z|=1/2}z^{-k-1}\sum_{n\le x}z^{\omega(n)}\dd z;
\]
the function $\Phi(z)=\frac{F(z)}{\Gamma(z)}$ is entire with $\Phi(z)=z+O(z^2)$.
\end{hint}

\section{Solutions}

\begin{solution}{ex:hankelint}
\idea Where $w^{-z}$ is single-valued, the loop can be closed into a circle and Hankel's formula becomes a residue computation.
\writeup For an integer $z$, $w^{-z}$ is single-valued, so the integrals along the two edges are equal and opposite and $\frac1{2\pi i}\int_{\mathcal H}w^{-z}e^w\dd w=\frac1{2\pi i}\oint_{|w|=1}w^{-z}e^w\dd w$. For $z=k\ge1$ this is the coefficient of $w^{k-1}$ in $e^w$, namely $\frac1{(k-1)!}=\frac1{\Gamma(k)}$: $1,1,\frac12$ for $k=1,2,3$. For $z=0$ or $-1$ the integrand $e^w$ or $we^w$ is entire and the integral vanishes, as does $\frac1\Gamma$ at $0$ and $-1$. For $z=\frac12$ the lemma gives $\frac1{\Gamma(1/2)}=\frac1{\sqrt\pi}$.
\end{solution}

\begin{solution}{ex:twoomega}
\idea Both proofs see the same constant: $\frac6{\pi^2}$ is the density of squarefree numbers, and $2^{\omega(n)}$ counts the squarefree divisors of $n$.
\writeup With $z=2$, $F(2)=\prod_p(1+\frac1p)(1-\frac1p)=\prod_p(1-\frac1{p^2})=\frac6{\pi^2}$ and $\Gamma(2)=1$, so Corollary~\ref{cor:omegasd} gives $x\log x(\frac6{\pi^2}+O(\frac1{\log x}))=\frac6{\pi^2}x\log x+O(x)$. Elementarily, $2^{\omega(n)}$ is the number of squarefree $d\mid n$, so $\sum_{n\le x}2^{\omega(n)}=\sum_{d\le x}\mu(d)^2\lfloor\frac xd\rfloor=x\sum_{d\le x}\frac{\mu(d)^2}d+O(x)$. By partial summation from $Q(y)=\sum_{d\le y}\mu(d)^2=\frac6{\pi^2}y+O(\sqrt y)$, $\sum_{d\le x}\frac{\mu(d)^2}d=\frac{Q(x)}x+\int_1^x\frac{Q(y)}{y^2}\dd y=\frac6{\pi^2}\log x+O(1)$.
\end{solution}

\begin{solution}{ex:minusone}
\idea At $z=-1$ the main term of the Selberg--Delange formula vanishes because $\frac1\Gamma$ has a zero at $-1$; only the error term remains.
\writeup Corollary~\ref{cor:omegasd} with $z=-1$, $R=1$: $\sum_{n\le x}(-1)^{\omega(n)}=x(\log x)^{-2}(\frac{F(-1)}{\Gamma(-1)}+O(\frac1{\log x}))$, and $\frac1{\Gamma(-1)}=0$, so the sum is $O(\frac x{(\log x)^3})$. In the language of the proof: $\zeta(s)^{-1}$ has no branch cut, the loop closes up, and the integrand $\zeta(s)^{-1}G(s)\frac{x^s}s$ is holomorphic at $s=1$. Delange's expansion of the Remark says more: every term contains $\frac1{\Gamma(-1-j)}=0$, so the sum is $O(\frac x{(\log x)^A})$ for every $A$. To do better one must use the zero-free region itself, as for the prime number theorem.
\end{solution}

\begin{solution}{ex:bigomegasd}
\idea Only the Euler factor changes; the method works as long as it stays bounded on $\sigma\ge\frac34$, which fails when $zp^{-s}$ can reach $1$.
\writeup Since $\Omega$ is completely additive, $\sum_nz^{\Omega(n)}n^{-s}=\prod_p(1-zp^{-s})^{-1}$ for $\sigma>1$ (absolute convergence: $|z|^{\Omega(n)}\le2^{\Omega(n)}\le\tau_2(n)$, so $K=2$ works). Then $G(s)=\prod_p\psi(p^{-s})$ with $\psi(u)=(1-u)^z(1-zu)^{-1}$. Here $\psi(0)=1$ and $\psi'(0)=-z+z=0$. For $|u|\le0.6$, $|zu|\le0.9$, so $\psi$ is holomorphic on the disc $|u|<\frac23$ and bounded on $|u|\le0.62$, and the maximum principle gives $|\psi(u)-1|\le C|u|^2$ for $|u|\le0.6$; the rest of Lemma~\ref{lem:eulerprod} applies verbatim, and Theorem~\ref{thm:sd} with $R=\frac32$ (so $K=2$, and $|z^{\Omega(n)}|\le\tau_2(n)$) gives the formula, with $G(1)=\prod_p(1-\frac1p)^z(1-\frac zp)^{-1}$. At $z=2$, the factor at $p=2$ is $(1-2^{1-s})^{-1}$, which has a pole at $s=1$: the integers $2^km$ contribute too much. In fact $\sum_{n\le x}2^{\Omega(n)}\ge\sum_{2^k\le x}2^k\#\{m\le x2^{-k}\text{ odd}\}$, about $\frac x2\log_2x$ already from this, and the true order is $x(\log x)^2$.
\end{solution}

\begin{solution}{ex:onlythree}
\idea Write the Euler product, compare it with $\zeta(s)/L(s,\chi_4)$, and take a square root.
\writeup The series is $F(s)=\prod_{p\equiv3\ (4)}(1-p^{-s})^{-1}$. Prime by prime, $\frac{\zeta(s)}{L(s,\chi_4)}$ has the factor $(1-2^{-s})^{-1}$ at $2$, $1$ at $p\equiv1$, and $\frac{1+p^{-s}}{1-p^{-s}}$ at $p\equiv3$. So $F^2=(1-2^{-s})\frac\zeta L\prod_{p\equiv3}(1-p^{-2s})^{-1}$, and comparing logarithms as in Corollary~\ref{cor:landau}, $F=\zeta^{1/2}G$ with $G=(1-2^{-s})^{1/2}L(s,\chi_4)^{-1/2}\prod_{p\equiv3}(1-p^{-2s})^{-1/2}$. Lemma~\ref{lem:logzeta} bounds $|L^{-1/2}|$ by $e^{A_2/2}(\log\tau)^{1/2}$, so Theorem~\ref{thm:sd} applies with $z=\frac12$, and the constant is $\frac{G(1)}{\sqrt\pi}=\frac1{\sqrt\pi}\cdot\frac1{\sqrt2}\cdot\frac2{\sqrt\pi}\prod_{p\equiv3}(1-p^{-2})^{-1/2}=\frac{\sqrt2}\pi\prod_{p\equiv3}(1-p^{-2})^{-1/2}$. The notebook finds the ratio of the count to $\frac{c_3x}{\sqrt{\log x}}$ to be $1.026$ at $x=10^7$.
\end{solution}

\begin{solution}{ex:landaupik}
\idea The Selberg--Delange formula holds uniformly on a circle around $z=0$, so Cauchy's formula extracts the coefficient of $z^k$; the main term is a Taylor coefficient of $e^{z\log\log x}\Phi(z)$.
\writeup Let $\lambda=\log\log x$ and $L=\log x$. Since $\sum_{n\le x}z^{\omega(n)}=\sum_k\pi_k(x)z^k$ is a polynomial,
\[
\pi_k(x)=\frac1{2\pi i}\oint_{|z|=1/2}\frac{\sum_{n\le x}z^{\omega(n)}}{z^{k+1}}\dd z=\frac xL\,\frac1{2\pi i}\oint_{|z|=1/2}\frac{e^{\lambda z}\Phi(z)}{z^{k+1}}\dd z+E,
\]
with $\Phi=\frac F\Gamma$ and, by Corollary~\ref{cor:omegasd} with $R=1$, $|E|\le\frac1{2\pi}\pi\,2^{k+1}\max_{|z|=1/2}|xL^{z-1}|\cdot O(\frac1L)=O_k(xL^{-3/2})$. The main integral is the coefficient of $z^k$ in $e^{\lambda z}\Phi(z)$. The function $\Phi$ is entire, with $\Phi(z)=\sum_{j\ge1}\phi_jz^j$, because $\frac1{\Gamma(z)}=z+O(z^2)$ and $F(0)=\prod_p(1-\frac1p)(1-\frac1p)^{-1}=1$; so $\phi_1=1$. The coefficient is $\sum_{j=1}^k\phi_j\frac{\lambda^{k-j}}{(k-j)!}=\frac{\lambda^{k-1}}{(k-1)!}(1+O_k(\frac1\lambda))$. Since $xL^{-3/2}=\frac xL\cdot L^{-1/2}$ is much smaller than $\frac xL\frac{\lambda^{k-2}}{(k-1)!}$, the claim follows. For $k=1$ this is the prime number theorem with $\pi_1(x)=\pi(x)+\#\{p^a\le x,a\ge2\}$, in the weak form $\frac x{\log x}(1+O(\frac1{\log\log x}))$.
\end{solution}

\section*{Sources and further reading}

\begin{itemize}
\item \emph{Landau and Ramanujan.} E.~Landau, ``Über die Einteilung der positiven ganzen Zahlen in vier Klassen nach der Mindestzahl der zu ihrer additiven Zusammensetzung erforderlichen Quadrate'', \emph{Arch.\ Math.\ Phys.}\ (3) 13 (1908), 305--312. B.~C.~Berndt and R.~A.~Rankin, \emph{Ramanujan: Letters and Commentary}, AMS/LMS, 1995, for the letter of 16 January 1913. G.~H.~Hardy, \emph{Ramanujan: Twelve Lectures on Subjects Suggested by His Life and Work}, Cambridge, 1940, pp.~9 and~63. D.~Shanks, ``The second-order term in the asymptotic expansion of $B(x)$'', \emph{Math.\ Comp.}\ 18 (1964), 75--86. P.~Moree, ``Counting numbers in multiplicative sets: Landau versus Ramanujan'', \emph{Mathematics Newsletter} 21 (2011), 73--81 (arXiv:1110.0708), for the story of Ramanujan's reply, of Stanley's computation and of Shanks's verdict; P.~Moree and J.~Cazaran, ``On a claim of Ramanujan in his first letter to Hardy'', \emph{Expo.\ Math.}\ 17 (1999), 289--312.
\item \emph{Sathe and Selberg.} L.~G.~Sathe, ``On a problem of Hardy on the distribution of integers having a given number of prime factors'', I--II, \emph{J.\ Indian Math.\ Soc.}\ 17 (1953), 63--141; III--IV, 18 (1954), 27--81. A.~Selberg, ``Note on a paper by L.~G.~Sathe'', \emph{J.\ Indian Math.\ Soc.}\ 18 (1954), 83--87; its opening paragraph calls Sathe's results very beautiful and his method complicated and involved.
\item \emph{Delange.} H.~Delange, ``Sur des formules dues à Atle Selberg'', \emph{Bull.\ Sci.\ Math.}\ (2) 83 (1959), 101--111; ``Sur des formules de Atle Selberg'', \emph{Acta Arith.}\ 19 (1971), 105--146.
\item \emph{Landau 1900.} E.~Landau, ``Sur quelques problèmes relatifs à la distribution des nombres premiers'', \emph{Bull.\ Soc.\ Math.\ France} 28 (1900), 25--38, for $\pi_k(x)$.
\item \emph{Hankel.} H.~Hankel, ``Die Euler'schen Integrale bei unbeschränkter Variabilität des Arguments'', \emph{Z.\ Math.\ Phys.}\ 9 (1864), 1--21.
\item \emph{Textbooks.} G.~Tenenbaum, \emph{Introduction to Analytic and Probabilistic Number Theory}, 3rd ed., AMS, 2015, Chapters~II.5--II.6, whose formulation of the method we follow. H.~L.~Montgomery and R.~C.~Vaughan, \emph{Multiplicative Number Theory~I}, Cambridge, 2007, Chapter~7. R.~de~la~Bretèche and G.~Tenenbaum, ``Remarks on the Selberg--Delange method'', \emph{Acta Arith.}\ 200 (2021), 349--369 (arXiv:2010.12929).
\end{itemize}
